\documentclass[11pt,a4paper]{article}

\usepackage[utf8]{inputenc}
\usepackage[T1]{fontenc}
\usepackage[main=swedish,provide=*]{babel}
\usepackage{amsmath,amssymb}
\usepackage{booktabs}
\usepackage[a4paper,margin=2.5cm]{geometry}
\usepackage[protrusion=true,expansion=false]{microtype}
\usepackage{parskip}

\DeclareMathOperator{\clip}{clip}

\begin{document}

\begin{center}
  {\LARGE\bfseries Den nya benchmarkingmodellen: en metodbeskrivning}\\[4pt]
  {\normalsize Juni 2026}
\end{center}

\begin{quote}\itshape\small
Status: Detta är en metodbeskrivning. Den återger vår operationalisering och tolkning av
Energimarknadsinspektionens inriktning för benchmarkingmodellen, given den information
myndigheten har offentliggjort fram till juni 2026. Vi skiljer genomgående mellan myndighetens
inriktning och våra arbetsantaganden. Där modellen vilar på arbetsantaganden, framför allt i
parametervärden och i hanteringen av data, är de utfall som följer av dem betingade på dessa
antaganden. Resultaten redovisas separat; detta dokument förklarar enbart metoden.
\end{quote}

\section{Syfte och läsanvisning}

Syftet med detta dokument är att förklara metoden bakom den nya benchmarkingmodellen för en
publik som känner till modellens beståndsdelar, alltså DEA-metoden, kapitalbasen och
kostnadsposterna, men inte nödvändigtvis hur de enskilda förändringarna är konstruerade eller
hur deras samlade effekt kan mätas.

Dokumentet är uppdelat i två delar som svarar på två olika frågor. Den första delen beskriver
förändringarna i själva modellen: hur kostnadsunderlaget byggs och hur effektiviseringskravet
beräknas. Den andra delen beskriver den metod vi använder för att mäta hur mycket var och en av
dessa förändringar bidrar till det samlade utfallet. Skillnaden är viktig: de första avsnitten
beskriver vad modellen gör, medan det sista beskriver hur vi räknar på vad modellen gör.

Genomgående skiljer vi på vad Energimarknadsinspektionen har angett som inriktning och vad som
är vår operationalisering. Myndigheten har offentliggjort en inriktning, medan flera detaljer
ännu inte är fastställda. Där vi har låst parametervärden eller gjort dataval för att kunna
räkna, markerar vi det, och de magnituder som följer är betingade på dessa val.

\section{Genomgående principer}

Två principer återkommer i flera av förändringarna, och vi formulerar dem en gång här.

\textbf{Normalisera i jämförelsen, ge incitament på de faktiska kostnaderna.} Flera av
förändringarna korrigerar för kostnadsskillnader som ett nätföretag inte kan påverka, exempelvis
elpriset i dess elområde eller förläggningsmiljön i dess nätområde. Korrigeringen görs i
benchmarkingen, så att företagen jämförs på mer lika villkor. Vid tillämpningen av incitamentet
görs däremot ingen motsvarande korrigering: incitamentet utgår från de faktiska, okorrigerade
kostnaderna, så att relativpriserna får genomslag i det slutliga utfallet i kronor. Samma logik
återkommer i värderingen av nätförluster och i korrigeringen för förläggningsmiljö.

\textbf{Relativ DEA innebär fördelningsmässiga effekter.} DEA-metoden mäter varje företag mot en
front av de mest effektiva företagen, och referensnivån bestäms ur samma fördelning. När en
kostnadspost förs in eller justeras förskjuts både det enskilda företagets position och fronten,
vilket påverkar samtliga företags relativa utfall. En förändrings effekt är därför
fördelningsmässig snarare än absolut: den omfördelar utfall mellan företag snarare än att skapa
eller ta bort kostnader i absoluta termer. Denna egenskap är avgörande både för hur
korrigeringarna ska tolkas och för hela dekompositionsmetoden i del II.

\section{Del I: Förändringarna i benchmarkingmodellen}

\subsection{Kostnadsunderlaget}

Den första gruppen av förändringar gäller hur det kostnadsunderlag som går in i DEA byggs.
Utgångspunkten är en TOTEX-ansats: samtliga kostnadsposter utom myndighetsavgifter ska ingå i en
och samma kostnadsvariabel. Skälet är en mer rättvisande jämförelse. Ett företag kan ha högre
kapitalkostnad men lägre kostnader för abonnemang till överliggande nät, medan ett annat har den
omvända sammansättningen. Bedöms företagen bara på en delmängd av kostnaderna jämförs de på sin
kostnadssammansättning snarare än på sin totala effektivitet. Myndighetsavgifter utesluts eftersom
de varken är påverkbara eller utbytbara, utan en bestämd summa som följer av antalet kunder och
deras spänningsnivåer.

\subsubsection{OPEX-variabeln}

OPEX-variabeln samlar de löpande kostnaderna. De påverkbara kostnaderna återanvänds från
baslinjen, så att den nya och den nuvarande modellen utgår från exakt samma påverkbara siffra och
därmed kan jämföras direkt. Till dessa läggs värderade nätförluster och ett urval icke-påverkbara
men utbytbara poster. Sammansättningen är
\[
  \mathrm{OPEX} = \text{påverkbara} \;+\; \text{värderade förluster}
  \;+\; \sum \text{utvalda icke-påverkbara} \;-\; \text{myndighetsavgifter},
\]
där de värderade förlusterna ges av
\[
  \text{värderade förluster} = \text{förlustfaktor} \times \text{gemensamt pris} \times
  \text{inmatad energi},
\]
summerade över prognosåren och medelvärdesbildade.

Nätförluster förs in i TOTEX, vilket samtidigt innebär att det fristående incitamentet för
nätförluster tas bort: en TOTEX-modell ger redan incitament att sänka förlusterna, eftersom en
lägre förlustkostnad sänker den totala kostnaden. I benchmarkingen värderas förlusterna till ett
gemensamt pris i stället för till varje företags faktiska förlustkostnad, eftersom den faktiska
kostnaden beror på elpriset i företagets elområde. Detta är principen ovan: prisskillnaden
neutraliseras i jämförelsen, medan incitamentet utgår från de faktiska kostnaderna. Myndighetens
inriktning tillåter att nätförluster, abonnemang eller båda värderas till ett gemensamt pris. I
nuläget gör modellen detta enbart för nätförluster, vilket följer av att tillräckligt detaljerad
data för abonnemang ännu saknas, inte av ett metodval.

Bland de icke-påverkbara posterna är abonnemang till överliggande och angränsande nät den enda
som inriktningen namnger uttryckligen. Den inkluderas trots att den är svår att påverka på kort
sikt, eftersom den är utbytbar mot kapitalkostnad. Ytterligare poster förs in under det generella
kriteriet att en kostnad ska vara i någon grad utbytbar eller påverkbar. Vilka dessa är följer av
en klassificering i datapipelinen snarare än av en uppräkning i inriktningstexten, vilket är värt
att vara tydlig med.

\subsubsection{Korrigeringen för förläggningsmiljö}

Det kostar olika mycket att anlägga elnät beroende på förläggningsmiljö, exempelvis i stadsmiljö
jämfört med på landsbygd, vilket ett företag normalt inte kan påverka utifrån sitt nätområde. Ett
företag med en dyrare förläggningsmiljö får en högre kapitalkostnad och faller därför ut sämre i
benchmarkingen på ett sätt som inte speglar dess effektivitet. Inriktningen är att justera ned
anskaffningsvärdet för jordkabel till en gemensam nivå, landsbygd normal, för de dyrare miljöerna.
Jordkabel ingår eftersom den utgör en stor del av kapitalbasen; kabelskåp utesluts som
försumbart; och en motsvarande korrigering görs för nätstationer.

Den metodiskt avgörande observationen är att förläggningsmiljöns premie redan ligger i prislistan.
Varje kabeltyps pris per kilometer är satt för den exakta kabeltypen i dess specifika
förläggningsmiljö, så en stadskabel har ett högre pris än samma kabeltyp på landsbygd.
Korrigeringen är därför inte ett mätproblem utan ett omprissättningsproblem: för att nivellera en
kabel ersätter man dess pris med landsbygd-normal-priset för samma kabeltyp. Premien behöver inte
skattas externt, utan är skillnaden mellan miljöns pris och landsbygd-normal-priset i samma lista.
Ur denna prisstruktur härleds ett schablonavdrag i procent per förläggningsmiljö, vilket är just
den form inriktningen efterfrågar för att slippa detaljerad rapportering per kabeltyp. Att
schablonen är tillräckligt träffsäker bekräftas av att den nära sammanfaller med en exakt
omprissättning per kabeltyp på sektornivå.

Två egenskaper i korrigeringen är lätta att läsa fel. För det första nivelleras hela
förläggningsmiljögradienten ned till landsbygd normal, inklusive landsbygd svår, så korrigeringen
är inte en lättnad enbart för stadsnät. För det andra justerar den bara nedåt: en kabel billigare
än landsbygd normal lämnas oförändrad. Korrigeringen kan därmed bara minska den mätta
kapitalbasen. Eftersom DEA är relativ är dess effekt på det enskilda företagets krav
fördelningsmässig, i enlighet med principen ovan.

\subsubsection{Sammanslagningen till en TOTEX-input}

I dagens modell ingår de påverkbara löpande kostnaderna och kapitalkostnaden som två separata
inputs i DEA. I den nya modellen slås de samman till en enda TOTEX-input. Skälet, enligt
inriktningen, är att dagens reglering ger incitament att inte investera i för dyra anläggningar men
inget incitament att välja den mest kostnadseffektiva lösningen, oavsett om den innebär löpande
kostnader eller investeringar. Med en sammanslagen kostnadsvariabel bedöms ett företag på sin
totala kostnad i stället för på fördelningen mellan löpande kostnader och kapital, vilket gör
jämförelsen lösningsneutral. En konsekvens är att den löpande respektive kapitalmässiga
effektiviteten inte längre kan bedömas var för sig; det är den totala kostnadseffektiviteten som
mäts.

\subsection{Kravmekaniken: det tvåsidiga effektiviseringskravet}

Den andra förändringen gäller hur effektiviseringskravet beräknas ur DEA-utfallet. Låt $E_i$
beteckna företag $i$:s effektivitet. I dagens modell mäts varje företag mot fronten ($E = 1$), och
potentialen, avståndet till fronten, är alltid icke-negativ. Ett golv i beräkningen gör att även
frontföretaget får ett avdrag på minst en viss nivå. Resultatet är att samtliga företag får ett
avdrag, och golvet fungerar i praktiken som ett generellt effektiviseringskrav som läggs på alla.

I den nya modellen flyttas referenspunkten från fronten till tredje kvartilen, alltså till
effektivitetsnivån $E_{75}$ vid den 75:e percentilen, beräknad exklusive outliers. Gapet definieras
som
\[
  \text{gap} = E_{75} - E_i,
\]
och kan nu bli negativt. Ett företag under tröskeln får ett avdrag, ett företag vid tröskeln får
full kostnadstäckning, och ett företag över tröskeln får ett tillägg. Eftersom tröskeln ligger vid
den 75:e percentilen är det per definition den översta fjärdedelen som får full täckning eller mer.
Valet av tredje kvartilen som referens är lånat från den brittiska regleringen. Tröskeln är
rörlig: den räknas om varje tillsynsperiod ur den aktuella fördelningen, så att den höjs när
branschen som helhet blir effektivare. I och med detta tas dagens golv och fasta outlierkrav bort,
och det finns inget separat generellt effektiviseringskrav i den nya modellen.

Själva beräkningskedjan från gap till årligt krav behåller i allt väsentligt sin form från dagens
modell. Det syns tydligast om de två formlerna ställs bredvid varandra, där $s$ är kundandelen,
$\clip$ kapar argumentet till det angivna intervallet, och faktorn $\tfrac{4}{8}$ skalar för att
potentialen realiseras över åtta år medan kravet sätts per fyraårig tillsynsperiod:
\begin{align*}
  \text{Nuvarande:} \quad
    & \text{krav} = \Bigl(1 + \clip\bigl(1 - E_i,\ \text{golv},\ \text{tak}\bigr)
      \cdot s \cdot \tfrac{4}{8}\Bigr)^{1/4} - 1, \\[2pt]
  \text{Ny:} \quad
    & \text{krav} = \Bigl(1 + \clip\bigl(E_{75} - E_i,\ -\text{tak},\ +\text{tak}\bigr)
      \cdot s \cdot \tfrac{4}{8}\Bigr)^{1/4} - 1.
\end{align*}
Det som har ändrats är alltså definitionen av gapet och att tecknet kan bli negativt. Tabell~\ref{tab:krav}
sammanfattar skillnaden.

\begin{table}[h]
  \centering
  \caption{Beräkningen av effektiviseringskravet i den nuvarande och den nya modellen.}
  \label{tab:krav}
  \begin{tabular}{lll}
    \toprule
                          & Nuvarande                & Ny \\
    \midrule
    Referenspunkt         & fronten ($E = 1$)        & tredje kvartilen ($E_{75}$) \\
    Tecken                & endast avdrag            & avdrag och tillägg \\
    Magnitudens drivare   & gap till fronten         & gap till $E_{75}$ \\
    Golv                  & ja, ger ett generellt krav & nej \\
    \bottomrule
  \end{tabular}
\end{table}

En strukturell egenskap är värd att lyfta: tvåsidigheten är skev. Effektiviteten är kapad vid sitt
högsta värde och tröskeln ligger nära toppen, så ett företag kan ligga bara en liten bit över
tröskeln men betydligt längre under den. Belöningssidan är därför strukturellt begränsad medan
avdragssidan kan bli stor, oavsett var ett symmetriskt tak sätts. Tvåsidigheten är alltså äkta men
i sak ojämn.

De konkreta parametervärdena i denna kedja, kundandelen, taket på gapet och realiseringsskalningen,
är ännu inte publicerade av myndigheten. Vi har låst dem som arbetsantaganden, valda för att spegla
dagens beräkningskedja, och redovisar dem i avsnittet om arbetsspecifikation. Takets nivå är
uttryckligen en av myndighetens öppna frågor.

\section{Del II: Att mäta förändringarnas samlade effekt}

De förändringar som beskrivits ovan verkar samtidigt, och vi vill kunna säga hur stor del av det
samlade utfallet som beror på var och en av dem, på ett sätt som summerar exakt till helheten. Den
uppenbara vägen, att slå av en förändring i taget och avläsa skillnaden, fungerar inte, och skälet
är principen om relativ DEA ovan. Eftersom varje komponent förskjuter fronten beror dess effekt på
vilka övriga komponenter som redan ingår. En komponents marginaleffekt är därför inte ett fast tal,
utan skiljer sig åt beroende på i vilken ordning komponenterna betraktas.

Låt $N$ vara mängden av komponenter, $n = |N|$, och låt $v(S)$ beteckna företagets signerade krav
när bara komponenterna i delmängden $S \subseteq N$ ingår. De två ytterlägena för en komponent $k$
är att ta bort den ur den fullständiga modellen och att lägga till den i den bara baslinjen:
\[
  \mathrm{LOO}(k) = v(N) - v(N \setminus \{k\}), \qquad
  \mathrm{AOI}(k) = v(\{k\}) - v(\emptyset).
\]
Eftersom DEA är icke-linjär sammanfaller dessa inte, och summan av de enskilda marginaleffekterna
träffar inte den faktiska totalen. Gapet mellan ytterlägena mäter hur starkt komponenten samverkar
med de övriga, och det är detta gap som motiverar en ordningsoberoende metod.

Shapley-värdet, hämtat från den kooperativa spelteorin, löser ordningsberoendet genom att beräkna
varje komponents marginaleffekt genomsnittad över samtliga ordningar i vilka komponenterna kan
läggas till modellen:
\[
  \varphi_k = \sum_{S \subseteq N \setminus \{k\}} w(|S|)\,
  \bigl[\, v(S \cup \{k\}) - v(S) \,\bigr], \qquad
  w(s) = \frac{s!\,(n - s - 1)!}{n!}.
\]
Det är den enda uppdelning som samtidigt summerar exakt till den totala förändringen, tilldelar en
komponent utan effekt exakt noll, behandlar utbytbara komponenter lika, och är additiv över
oberoende delar. Den första av dessa egenskaper, som är hela poängen med uppdelningen, är
identiteten
\[
  \sum_{k \in N} \varphi_k = v(N) - v(\emptyset).
\]
Resultatet är en exakt och rättvis bokföring av vad modellen gör.

Förändringarna i del I är just de komponenter som dekompositionen fördelar, men de är inte av samma
slag, och vi delar därför upp dekompositionen i två skikt. De beräkningsmässiga förändringarna,
sammanslagningen till en TOTEX-input och övergången till det tvåsidiga kravet, utgör ett yttre
skikt: de bestämmer hur kravet beräknas. Kostnadsunderlagets enskilda poster utgör ett inre skikt:
de bestämmer vilka kostnader som ingår. Det yttre skiktet fördelar steget från dagens krav till den
nya modellens baslinje, och det inre skiktet fördelar därefter kravet mellan komponenterna, givet
den nya beräkningsregimen. Båda skikten är exakta och summerar tillsammans till hela förändringen.
Uppdelningen är medveten: komponenternas bidrag redovisas givet den modell som faktiskt gäller,
inte genomsnittat över en avskaffad beräkningsregim.

Två förbehåll följer av principen om relativ DEA. För det första är bidragen fördelningsmässiga,
inte absoluta: ett bidrag som sänker ett företags krav motsvaras i regel av en höjning hos andra, så
en komponents bidrag över företagen summerar nära noll. När vi säger att en komponent gynnar ett
företag är det ett påstående om dess läge i förhållande till jämförelsebolagen, inte om kostnader i
absoluta termer. För det andra kan en komponent vara liten i genomsnitt men stor i omfördelning:
genomsnittet säger hur den påverkar sektorns samlade krav, medan spridningen säger hur den flyttar
om de enskilda företagen. Dessa två frågor besvaras därför var för sig i resultatredovisningen.
Slutligen är Shapley-värdet en rättvis bokföring av vad modellen gör, givet modellvalen, inte en
kausal kontrafaktisk utsaga.

\section{Arbetsspecifikation och förbehåll}

Följande punkter är våra arbetsantaganden eller följer av tillgänglig data, och bör läsas som sådana
tills myndighetens specifikation publiceras. Samtliga magnituder som redovisas separat är betingade
på dem.

\textit{Det tvåsidiga kravet.}
\begin{enumerate}
  \item Kundandelen, taket på gapet och realiseringsskalningen är inte publicerade av myndigheten.
    Vi har låst dem för att spegla dagens beräkningskedja. Takets nivå är en öppen regleringsfråga.
  \item Att magnituden bestäms av det kardinala gapet, och att golv och fast outlierkrav tas bort,
    är vår tolkning av en funktionsform som myndigheten inte har fastställt.
\end{enumerate}

\textit{Förläggningsmiljökorrigeringen.}
\begin{enumerate}\setcounter{enumi}{2}
  \item De konkreta procentsatserna är vår härledning ur prisstrukturen; myndighetens officiella
    schablonavdrag är ännu inte publicerat.
  \item Premien härleds ur normvärdeslistan, medan den kommande metoden riktar sig mot
    anskaffningsvärde, som saknas i nära all tillgänglig data. Härledningen vilar därmed på att
    prisrelationen mellan förläggningsmiljöer är densamma i anskaffningsvärde som i normvärde.
    Antagandet är rimligt men inte verifierbart ur denna data.
\end{enumerate}

\textit{OPEX-variabeln.}
\begin{enumerate}\setcounter{enumi}{4}
  \item Att i nuläget bara värdera nätförluster, och inte abonnemang, till ett gemensamt pris följer
    av en databegränsning, inte av ett metodval.
  \item Urvalet av icke-påverkbara poster utöver abonnemang vilar på det generella
    utbytbarhetskriteriet och är klassificerat i datapipelinen.
  \item Blandningen av perioder och valet av prognosperiod är kända modellförenklingar som speglar
    dagens modell.
\end{enumerate}

\section{Sammanfattning}

\begin{enumerate}
  \item Kostnadsunderlaget byggs som en TOTEX: samtliga poster utom myndighetsavgifter samlas i en
    kostnadsvariabel, så att företagen jämförs på total kostnad i stället för på
    kostnadssammansättning, och de två tidigare inputerna slås samman till en.
  \item Nätförluster värderas till ett gemensamt pris och förläggningsmiljön nivelleras till en
    gemensam nivå, båda enligt principen att normalisera i jämförelsen men ge incitament på de
    faktiska kostnaderna.
  \item Effektiviseringskravet blir tvåsidigt: referensen flyttas från fronten till tredje
    kvartilen, kravet kan bli negativt, den översta fjärdedelen får full täckning eller mer, och
    dagens golv tas bort. Tvåsidigheten är strukturellt skev.
  \item För att mäta varje förändrings bidrag till det samlade utfallet används Shapley-värdet, som
    ger en exakt och ordningsoberoende uppdelning. Bidragen är fördelningsmässiga och betingade på
    modellvalen, inte absoluta eller kausala.
  \item Parametervärden och vissa dataval är arbetsantaganden, inte myndighetens beslut, och de
    magnituder som redovisas separat är betingade på dem.
\end{enumerate}

\end{document}
